\par\noindent\textbf{(1)} 工质每循环状态回复，ΔS\_sys=0，热库熵变化为−Qh/Th+Qc/Tc。总熵产生S\_gen=Qc/Tc−Qh/Th≥0，因此放热必须满足Qc≥(Tc/Th)Qh。式中不可把热机内部可逆误写成热库总熵为负。
\par\noindent\textbf{(2)} 效率η=W/Qh=1−Qc/Qh，将上式代入得η≤1−Tc/Th=η\_C。更精确有η\_C−η=(Tc/Qh)S\_gen，说明任意有限不可逆熵产都使效率严格低于Carnot极限，而不仅是口头的第二定律描述。
\par\noindent\textbf{(3)} 固定Qh且对同一两热库工作，理想可得W\_rev=Qh(1−Tc/Th)，实际W=W\_rev−TcS\_gen。损失功W\_lost=TcS\_gen即Gouy–Stodola关系，Tc为作功环境温度；若环境温度不是固定冷库，须重新定义火用参考态。
\par\noindent\textbf{(4)} 若热源是有限热容量物体，放热过程中温度Th(T)变化，热源熵变应积分dS=δQ\_rev/T，而非用总热量除以某一个初温。热机最大可逆功由系统总熵守恒加能量守恒共同确定；简单η\_C公式只适用两个无限等温热库。
